% Supporting information template.

\documentclass[aps,prl,preprint,preprintnumbers,amsmath,amssymb,superscriptaddress]{revtex4-2}

\usepackage{amsmath}
\usepackage{amssymb}
\usepackage{amsthm}
\usepackage{booktabs}
\usepackage{braket}
\usepackage{chemformula}
\usepackage{float}
\usepackage{graphicx}
\usepackage{mathrsfs}
\usepackage{multirow}
\usepackage{physics}
\usepackage{tabularx}
\usepackage{threeparttable}
\usepackage{xcolor}

\renewcommand{\thetable}{S\arabic{table}}
\renewcommand{\thefigure}{S\arabic{figure}}

\begin{document}

\title{Supporting Information for ``Manuscript Title''}

\author{First Author}
\affiliation{Department, Institution, City, Country}

\author{Second Author}
\affiliation{Department, Institution, City, Country}

\maketitle

\section{Additional Methods}

Add supplementary methods, derivations, or computational details here.

\section{Additional Results}

Add supplementary figures, tables, and discussion here.

\begin{figure}[H]
    \centering
    % \includegraphics[width=0.8\textwidth]{example-supporting-figure.pdf}
    \caption{Replace this caption with a concise description of the supporting figure.}
    \label{fig:supporting-example}
\end{figure}

\begin{table}[H]
    \centering
    \caption{Replace this caption with a concise description of the supporting table.}
    \label{tab:supporting-example}
    \begin{tabular}{lll}
        \toprule
        Column 1 & Column 2 & Column 3 \\
        \midrule
        Entry & Entry & Entry \\
        \bottomrule
    \end{tabular}
\end{table}

\end{document}
